% ============================================================================
%  Subdocumento 3. Contenido. Se usa igual en el documento único y en el
%  PDF propio del subdocumento. Los formularios que le asigna el T-21 se
%  agregan solos al final (datos en formularios/ de esta carpeta).
% ============================================================================

\begin{resumenApertura}
\marcador{Qué resuelve este subdocumento, en cuatro o cinco líneas}
\begin{recibe}
  \item \marcador{Qué recibe el mandante}
\end{recibe}
\end{resumenApertura}

\section{Esquema de la solución}\label{sec:3-esquema-de-la-solucion}

\marcador{Modelo conceptual con diagrama, componentes y trazabilidad a los requerimientos}

\section{Alcance de la Etapa 1 y de la Etapa 2}\label{sec:3-alcance-de-la-etapa-1-y-de-la-et}

\marcador{Separación explícita, criterios de asignación, dependencias e hitos externos}

\section{Exclusiones, supuestos y restricciones}\label{sec:3-exclusiones-supuestos-y-restricc}

\marcador{Exclusiones del Caso, supuestos del alcance y restricciones}

\section{Catálogo de requerimientos}\label{sec:3-catalogo-de-requerimientos}

\marcador{Resumen y análisis del catálogo. El detalle va en el Formulario T-12}

\section{Estrategia con los grupos de interés}\label{sec:3-estrategia-con-los-grupos-de-int}

\marcador{Plan de adhesión y apoyo de los involucrados clave}

\section{Criterios de aceptación del alcance}\label{sec:3-criterios-de-aceptacion-del-alca}

\marcador{Los criterios del Caso con meta, hito y medición}
